\documentclass[a4paper,12pt]{report}
\usepackage[utf8]{inputenc}
\usepackage[T1]{fontenc}
\usepackage[french]{babel}
\usepackage{array,enumitem}
\usepackage[breaklinks=true,hidelinks,hyperfootnotes=false,linkcolor=black,citecolor=black]{hyperref}
\usepackage{xurl}
\usepackage{titlesec}
\usepackage{graphicx}
\usepackage{geometry}
\usepackage{csquotes}
\usepackage{longtable}
\geometry{a4paper, top=0.8in, bottom=1in, left=1in, right=1in}
\usepackage{setspace}
\usepackage{mathptmx}
\usepackage{xcolor}
\usepackage{amsmath}
\definecolor{customblue}{RGB}{31, 77, 121}
\definecolor{teal100}{rgb}{0.852, 0.948, 0.928}
\definecolor{headerblue}{RGB}{31, 77, 121}
\usepackage{float}
\usepackage{caption}
\captionsetup{font=scriptsize}
\usepackage[table,xcdraw]{xcolor}
\usepackage{colortbl}
\usepackage{booktabs}

\titleformat{\chapter}[hang]
  {\normalfont\bfseries\Large\color{customblue}}
  {Chapitre \thechapter\ —}
  {0.5em}{}
  [\vspace{0.3ex}\rule{\textwidth}{0.8pt}]
\titlespacing*{\chapter}{0pt}{10pt}{10pt}

\titleformat{\section}{\normalfont\bfseries\large}{}{0em}{}
\titleformat{\subsection}{\normalfont\bfseries\normalsize}{}{0em}{}

\onehalfspacing

% ─── DOCUMENT ────────────────────────────────────────────────────────────────
\begin{document}
\pagenumbering{gobble}

% ─── PAGE DE GARDE ───────────────────────────────────────────────────────────
\begin{titlepage}
\centering
\input{page_garde}

\end{titlepage}
\clearpage

% ─── TABLE DES MATIÈRES ──────────────────────────────────────────────────────
\pagenumbering{roman}
\tableofcontents
\clearpage

% ─── INTRODUCTION ────────────────────────────────────────────────────────────
\pagenumbering{arabic}
\setcounter{page}{1}

\chapter*{Introduction générale}
\addcontentsline{toc}{chapter}{Introduction générale}

La théorie des graphes constitue l'un des fondements les plus puissants de l'informatique moderne.
Elle intervient dans des domaines aussi variés que la navigation GPS, les réseaux de communication,
l'intelligence artificielle ou encore l'optimisation logistique.
Pourtant, l'apprentissage de ses algorithmes reste souvent difficile en raison de leur caractère
abstrait et de la complexité des structures manipulées.

C'est dans ce contexte que s'inscrit le projet \textit{Graph Lab} : une plateforme interactive
permettant de créer, manipuler et visualiser des graphes, tout en exécutant pas à pas les principaux
algorithmes classiques. L'objectif est de rendre ces concepts accessibles grâce à une expérience
visuelle et interactive.

Ce rapport présente l'ensemble du travail réalisé par l'équipe, depuis la conception du système
jusqu'à la livraison du produit final, en passant par la méthodologie adoptée et la répartition
des tâches.


\clearpage

% ═══════════════════════════════════════════════════════════════════════════
\chapter{Présentation du projet}
% ═══════════════════════════════════════════════════════════════════════════

\section{Problématique}

Malgré leur importance en informatique et en mathématiques discrètes, les algorithmes de graphes restent souvent difficiles à appréhender. Cette difficulté s'explique par plusieurs facteurs :

\begin{itemize}
    \item Leur nature abstraite et théorique.
    \item Le manque d'outils interactifs permettant leur manipulation dynamique.
    \item Une compréhension souvent limitée aux explications statiques (cours, schémas).
\end{itemize}

Ces limitations rendent l'apprentissage et l'expérimentation des algorithmes de graphes peu intuitifs.

\section{Objectifs}

Dans ce cadre, ce projet vise à développer une plateforme interactive dédiée à l'expérimentation des algorithmes de graphes. Les principaux objectifs sont :

\begin{itemize}
    \item Permettre la création et la manipulation de graphes de manière simple et intuitive.
    \item Exécuter et observer le fonctionnement des algorithmes classiques de graphes.
    \item Améliorer l'apprentissage par l'interaction et la visualisation dynamique.
\end{itemize}


\section{Fonctionnalités de la plateforme}

La plateforme développée permet :

\begin{itemize}
    \item La création de graphes orientés ou non orientés, pondérés ou non pondérés.
    \item L'ajout et la suppression de sommets et d'arêtes via un éditeur visuel SVG ou une saisie matricielle.
    \item Le chargement ou la génération de graphes pour les tests.
    \item L'exécution des algorithmes suivants :
    \begin{itemize}
        \item Problème de coloration de graphes (greedy coloring)
        \item Recherche de composantes connexes et fortement connexes
        \item Plus courts chemins (Dijkstra, Bellman-Ford)
        \item Arbres couvrants minimaux (Prim, Kruskal)
        \item Flot maximum (Edmonds-Karp)
        \item Recherche de circuits eulériens (Hierholzer)
    \end{itemize}
    \item La visualisation graphique des résultats ainsi que des étapes d'exécution des algorithmes (Algorithm Cinema).
    \item Un mini-langage de requêtes interactif (\texttt{path}, \texttt{neighbors}, \texttt{degree}, \texttt{components}).
    \item La gestion des cas limites (graphes vides, présence de cycles négatifs, poids négatifs, graphes non eulériens, etc.).
\end{itemize}

\section{Technologies utilisées}

\begin{itemize}
    \item \textbf{React.js} — interface utilisateur et gestion des interactions
    \item \textbf{TypeScript} — structuration et sécurité statique du code
    \item \textbf{D3.js} — visualisation dynamique des graphes sur canvas SVG
    \item \textbf{Tailwind CSS} — design et mise en forme responsive
    \item \textbf{Vite} — environnement de build et serveur de développement
\end{itemize}

\section{Architecture globale du système}

La plateforme repose sur une architecture logique unifiée, sans séparation stricte entre
frontend et backend. Elle est organisée en trois modules principaux fortement intégrés,
permettant une exécution fluide et une visualisation en temps réel des résultats.

\begin{itemize}
    \item \textbf{Module Interface utilisateur} : gestion des interactions, création et
          modification des graphes via un éditeur visuel SVG ou une saisie matricielle.
    \item \textbf{Module Logique des graphes (\texttt{GraphUI})} : gestion des structures
          de données (sommets, arêtes, poids, orientations) et du contrat de données
          partagé entre les modules, géré par un reducer TypeScript centralisé.
    \item \textbf{Module Algorithmes (\texttt{algorithmCinema})} : implémentation et exécution
          des différents algorithmes de graphes avec génération des snapshots de visualisation.
\end{itemize}

\clearpage


% ═══════════════════════════════════════════════════════════════════════════
\chapter{Organisation de l'équipe et méthode de travail}
% ═══════════════════════════════════════════════════════════════════════════

\section{Composition et rôles}

L'équipe de onze membres est organisée autour de trois pôles complémentaires, avec une cheffe
de projet, des développeurs algorithmes, et deux équipes transversales couvrant l'interface et
les tests.

\textbf{Chef de projet} : Coordination générale, distribution des tâches, orchestration des
sprints, documentation et intégration des modules.

\textbf{Développeurs algorithmes} : Chaque développeur est responsable de l'implémentation
complète d'un ou plusieurs algorithmes en TypeScript, incluant la génération des snapshots
de visualisation.

\textbf{Équipe front-end} : Rayen Mestiri et Yassine Drira forment le binôme principal en
charge de l'interface utilisateur. Ils ont également contribué aux tests de leurs composants.

\textbf{Équipe test} : Youssef Fathallah et Idriss Labibi constituent le binôme dédié à la
validation. Ils ont également contribué à l'interface en développant des composants d'affichage
des résultats. Les deux équipes ont donc collaboré sur l'interface et les tests.

\section{Répartition détaillée des tâches}

\renewcommand{\arraystretch}{1.4}
\begin{longtable}{>{\bfseries}p{3.6cm} p{2.8cm} p{7.2cm}}
\rowcolor{headerblue}
\color{white}\textbf{Membre} & \color{white}\textbf{Rôle} & \color{white}\textbf{Tâches réalisées} \\
\hline
\endfirsthead
\rowcolor{headerblue}
\color{white}\textbf{Membre} & \color{white}\textbf{Rôle} & \color{white}\textbf{Tâches réalisées} \\
\hline
\endhead

\rowcolor{teal100}
Ben Boubaker Tasnime & Chef de projet &
Création du dépôt GitHub, organisation des daily standups, distribution et suivi des tâches,
intégration des modules, rédaction de la documentation et du contrat \texttt{GraphUI}. \\

Omar Abdelkader & Dév. Algo &
Algorithme de coloration gloutonne (\textit{greedy coloring}) : implémentation TypeScript, calcul du nombre chromatique, snapshots de colorisation, tests sur graphes bipartis et planaires. \\

\rowcolor{teal100}
Rihem Amri & Dév. Algo &
Composantes connexes (BFS/Union-Find) et fortement connexes (Kosaraju et Tarjan) : implémentation, snapshots de découverte progressive, tests sur graphes orientés et non orientés. \\

Nada Ben Taher & Dév. Algo &
Dijkstra : implémentation avec file de priorité (min-heap), snapshots de relaxation des arêtes, gestion des graphes sans poids négatifs, tests de cohérence. \\

\rowcolor{teal100}
Nour El Houda Hasnaoui & Dév. Algo &
Bellman-Ford : implémentation avec gestion des poids négatifs, détection des cycles négatifs, snapshots itération par itération, alerte visuelle en cas de cycle négatif. \\

Guerriche Oussema & Dév. Algo &
Algorithme de Prim (ACM) et flot maximum (Edmonds-Karp) : implémentation, snapshots d'évolution du flot et de l'arbre couvrant, gestion du graphe résiduel. \\

\rowcolor{teal100}
Rayen Mestiri & \textbf{Équipe front} / Dév. Algo &
\textbf{Interface :} barre d'outils, contrôles de mode (visuel / matrice), panneau de contrat \texttt{GraphUI} en temps réel, saisie manuelle et import de graphes. \textbf{Algo :} Kruskal (ACM) — tri des arêtes, Union-Find avec compression de chemin et union par rang, snapshots accepté/rejeté. \textbf{Tests :} validation des composants front développés. \\

AmenAllah Hajji & Dév. Algo &
Hierholzer (circuits eulériens) : vérification des conditions d'Euléricité (degrés pairs / équilibre in/out), construction du circuit par combinaison de sous-circuits, snapshots pas à pas, messages d'erreur pour graphes non eulériens. \\

\rowcolor{teal100}
Drira Yassine & \textbf{Équipe front} &
Menu interactif principal, sélection du type de graphe et de l'algorithme, panneau de résultats en temps réel, intégration avec le reducer \texttt{GraphUI}, mini-langage de requêtes en canvas (\texttt{path}, \texttt{neighbors}, \texttt{degree}, \texttt{components}). \textbf{Tests :} scénarios nominaux sur l'éditeur visuel et la saisie matricielle. \\

Labibi Idris & \textbf{Équipe test} / Front &
\textbf{Tests :} validation sur arbres, chemins linéaires $P_n$, cycles courts $C_3$--$C_5$ et graphes complets $K_n$ ($n \leq 6$). \textbf{Front :} composants d'affichage des résultats algorithmiques et du panneau de lecture des snapshots. \\

\rowcolor{teal100}
Fathallah Youssef & \textbf{Équipe test} / Front &
\textbf{Tests :} cas limites — graphe vide, graphe non connexe, poids négatifs, cycles négatifs, graphe sans circuit eulérien, boucles et arêtes multiples. \textbf{Front :} composants d'affichage des alertes et messages d'erreur dans l'interface. \\

\end{longtable}

\section{Méthode de travail}

L'équipe a adopté une organisation agile avec sprints hebdomadaires, chacun correspondant à un livrable précis. Les \textbf{daily standups} (10--15 min) ont permis de synchroniser les avancées, détecter les dépendances inter-modules (ex. : l'interface ne peut afficher les étapes qu'après la génération des snapshots par les modules algo) et résoudre rapidement les blocages.

La collaboration entre l'équipe front (Rayen, Yassine) et l'équipe test (Youssef, Idriss) a été particulièrement productive : les testeurs ont identifié des bugs d'affichage et contribué à des composants UI, tandis que les développeurs front ont participé à la rédaction de cas de test pour leurs propres composants.

La méthodologie suit cinq phases séquentielles :

\begin{enumerate}[leftmargin=1.5cm]
  \item \textbf{Conception} : définition du contrat \texttt{GraphUI}, interfaces TypeScript, architecture des modules.
  \item \textbf{Développement parallèle} : chaque développeur travaille indépendamment en respectant le contrat commun.
  \item \textbf{Intégration} : fusion des branches par la cheffe de projet, vérification de la cohérence.
  \item \textbf{Tests} : validation sur cas nominaux et cas limites par les deux équipes.
  \item \textbf{Livraison} : documentation finale et démo de la plateforme.
\end{enumerate}

\textbf{Outils de collaboration :} GitHub (versionnage, branches par fonctionnalité, pull requests), messagerie instantanée (échanges rapides), npm/Node.js/Vite (build partagé).

\clearpage


% ═══════════════════════════════════════════════════════════════════════════
\chapter{Travail réalisé et produit livrable}

 

\section{Conception du système}

\subsection{Contrat de données : la classe \texttt{GraphUI}}

Le cœur de l'architecture repose sur un contrat de données unique, la classe \texttt{GraphUI},
qui constitue la source de vérité partagée entre tous les modules. Sa définition TypeScript est
la suivante :

\begin{verbatim}
class GraphUI {
  nodes: number[]
  edges: Array<{ from: number; to: number; weight: number }>
  directed: boolean
  weighted: boolean
  positions?: Record<number, { x: number; y: number }>
}
\end{verbatim}

Ce contrat garantit que l'interface produit exactement le format attendu par la couche
algorithmique, sans conversion intermédiaire. La prévisualisation JSON de ce contrat est
affichée en temps réel dans le panneau droit de l'application.

\subsection{Choix des structures de données}

Les choix de structures de données ont été guidés par deux impératifs : la performance
algorithmique et la compatibilité avec la visualisation en temps réel.

\begin{itemize}
    \item \textbf{Listes d'adjacence} : utilisées pour la majorité des algorithmes de parcours
          (BFS, DFS, composantes connexes, Euler) car elles offrent une complexité spatiale
          proportionnelle au nombre d'arêtes.
    \item \textbf{Matrice d'adjacence} : disponible via le module de saisie matricielle,
          particulièrement adaptée aux graphes denses et à la détection rapide d'arêtes.
    \item \textbf{Union-Find (structure disjointe)} : utilisée par Kruskal et la détection
          de composantes connexes pour des opérations quasi-linéaires.
    \item \textbf{File de priorité (min-heap)} : utilisée par Dijkstra et Prim pour
          l'extraction efficace du sommet de poids minimal.
    \item \textbf{Reducer TypeScript centralisé} : l'état canonique du graphe est géré par
          un reducer Redux-like, garantissant l'immuabilité et la traçabilité des mutations
          avec un historique complet (undo/redo/jump).
\end{itemize}

 
% ─────────────────────────────────────────────────────────────────────────
\section{Interface utilisateur}
% ─────────────────────────────────────────────────────────────────────────
 
\subsection{Éditeur visuel (canvas SVG — D3.js)}
 
L'éditeur SVG, développé principalement par l'équipe front, gère les interactions suivantes :
ajout de nœuds par clic, déplacement par glisser-déposer, création d'arêtes entre deux sommets,
suppression, et affichage de \textit{diff ghosts} (états précédents en transparence) lors du
time-travel. Un scrubber de timeline permet de naviguer dans l'historique des mutations
(\texttt{UNDO} / \texttt{REDO} / \texttt{JUMP\_TO}).
 
\subsection{Saisie matricielle}
 
Le module matriciel permet de redimensionner le nombre de sommets, d'éditer les cellules
de la matrice d'adjacence et d'appliquer la matrice au graphe, avec synchronisation
automatique de l'éditeur visuel. Les deux modes partagent le même état \texttt{GraphUI}.
 
\subsection{Algorithm Cinema et mini-langage de requêtes}

Le module \textit{Algorithm Cinema} (fichier \texttt{algorithmCinema.ts}) permet une visualisation
pédagogique des algorithmes :
\begin{itemize}
    \item \textbf{Lecture / Pause / Étape suivante / Rembobinage} pour naviguer dans les snapshots.
    \item \textbf{Scrubber} pour sauter directement à n'importe quelle étape de l'exécution.
    \item Mise en évidence visuelle des arêtes et sommets actifs à chaque étape.
\end{itemize}

En complément, un \textbf{mini-langage de requêtes} est disponible directement en canvas :
\texttt{path 1 5} (chemin entre deux sommets), \texttt{neighbors 3} (voisins d'un nœud),
\texttt{degree 4} (degré d'un sommet), \texttt{components} (liste des composantes connexes).

 

% ─────────────────────────────────────────────────────────────────────────
\section{Implémentation des algorithmes}
% ─────────────────────────────────────────────────────────────────────────

Un graphe est défini par $G = (V, E)$, où $V$ représente l'ensemble des sommets et $E$ l'ensemble des arêtes reliant ces sommets.

% ─────────────────────────────────────────────────────────────────────────
\subsection{Plus courts chemins}

\textbf{Dijkstra} — \textit{Nada Ben Taher.}

Calcule les plus courts chemins depuis une source dans un graphe à poids \emph{strictement positifs}.
L'algorithme maintient une \textbf{file de priorité (min-heap)} et extrait à chaque étape le sommet
de distance minimale provisoire. Pour chaque voisin de ce sommet, il effectue une \textbf{relaxation} :
si le chemin passant par le sommet courant est plus court que la distance connue, la distance est mise
à jour. Ce processus se répète jusqu'à ce que tous les sommets atteignables soient traités. Chaque
relaxation génère un snapshot pour la visualisation pas à pas. Complexité : $O((V+E)\log V)$.

\medskip
\textbf{Bellman-Ford} — \textit{Nour El Houda Hasnaoui.}

Calcule les plus courts chemins dans les graphes \emph{avec poids négatifs}. Il effectue $|V|-1$
itérations de mise à jour sur la totalité des arêtes. Une $|V|$-ème itération permet de détecter
l'existence d'un \textbf{cycle de poids négatif} : si une distance peut encore être améliorée, un
cycle négatif est présent. Un snapshot est généré à chaque itération, avec une alerte visuelle
indiquant l'arête incriminée en cas de cycle. Complexité : $O(V \cdot E)$.
 
\subsection{Arbres couvrants minimaux}

\textbf{Prim} — \textit{Guerriche Oussema.}

Construit progressivement un arbre couvrant minimal à partir d'un sommet initial arbitraire.
L'algorithme maintient une \textbf{file de priorité} contenant les arêtes frontières reliant
l'arbre courant aux sommets non encore visités. À chaque étape, l'arête de poids minimal est
sélectionnée et son sommet cible est intégré à l'arbre. Un snapshot est généré après chaque
ajout. Complexité : $O((V+E)\log V)$.

\medskip
\textbf{Kruskal} — \textit{Rayen Mestiri.}

Commence par trier toutes les arêtes par poids croissant. Chaque arête est ensuite examinée
à l'aide de la structure \textbf{Union-Find} (avec compression de chemin et union par rang).
Si les deux extrémités appartiennent à des composantes différentes, l'arête est ajoutée à
l'arbre couvrant et un snapshot \textit{accepté} est généré. Sinon, l'arête est rejetée
(elle formerait un cycle) et un snapshot \textit{rejeté} est affiché. Complexité : $O(E \log E)$.

% ─────────────────────────────────────────────────────────────────────────

\subsection{Flot maximum}

\textbf{Edmonds-Karp} — \textit{Guerriche Oussema.}

Variante de Ford-Fulkerson reposant sur des recherches successives de \textbf{chemins augmentants}
dans le graphe résiduel à l'aide d'un BFS (garantissant le plus court chemin en nombre d'arêtes).
À chaque itération, le flot est augmenté selon la capacité résiduelle minimale du chemin trouvé,
puis les capacités résiduelles sont mises à jour (sens aller et retour). Un snapshot visualise
l'état du graphe résiduel et la valeur courante du flot. L'algorithme se termine lorsqu'aucun
chemin augmentant n'existe entre la source et le puits (coupe minimale atteinte).
Complexité : $O(VE^2)$.

% ─────────────────────────────────────────────────────────────────────────

\subsection{Composantes connexes et fortement connexes}

\textbf{Composantes connexes} — \textit{Rihem Amri.}

Sur un graphe non orienté, un parcours BFS ou DFS est lancé depuis chaque sommet non encore visité.
Chaque parcours identifie une composante distincte, colorée différemment dans la visualisation.
Une alternative utilise la structure \textbf{Union-Find} pour fusionner progressivement les sommets
reliés par une arête. La requête \texttt{components} du mini-langage exploite ce module.

\medskip
\textbf{Composantes fortement connexes} — \textit{Rihem Amri.}

Deux approches sont implémentées sur graphe orienté. \textit{Kosaraju} effectue d'abord un DFS
complet pour ordonner les sommets selon leur temps de fin, puis relance un DFS sur le graphe
transposé dans l'ordre inverse. \textit{Tarjan} repose sur un DFS unique utilisant une pile
et des indices \textit{low-link} pour identifier les racines de composantes fortement connexes
sans transposition du graphe. Dans les deux cas, des snapshots illustrent la découverte progressive.
Complexité : $O(V+E)$.

% ─────────────────────────────────────────────────────────────────────────

\subsection{Coloration de graphes}

\textbf{Greedy coloring} — \textit{Omar Abdelkader.}

Les sommets sont traités dans un ordre fixe (ordre d'insertion). Pour chacun, la plus petite
couleur disponible — c'est-à-dire non utilisée par ses voisins déjà colorés — est attribuée.
Le résultat fournit une coloration valide (aucune arête ne relie deux sommets de même couleur)
et le \textbf{nombre chromatique expérimental} est affiché dans l'interface. L'algorithme est
validé sur des graphes bipartis (2 couleurs suffisent), des graphes planaires (4 couleurs au
plus) et des graphes complets $K_n$ ($n$ couleurs nécessaires). Complexité : $O(V + E)$.

% ─────────────────────────────────────────────────────────────────────────

\subsection{Circuits eulériens}

\textbf{Hierholzer} — \textit{AmenAllah Hajji.}

Avant exécution, le graphe est vérifié pour les \textbf{conditions d'Euléricité} :
dans un graphe non orienté, tous les sommets doivent avoir un degré pair ;
dans un graphe orienté, il faut $\deg^+(v) = \deg^-(v)$ pour chaque sommet.
Si la condition n'est pas satisfaite, les sommets fautifs sont listés et l'exécution est stoppée
avec un message explicatif.

Sinon, l'algorithme construit le circuit en combinant des sous-circuits successifs.
Il démarre depuis un sommet source, parcourt un premier cycle en explorant des arêtes non visitées,
puis insère progressivement d'autres cycles aux sommets disposant encore d'arêtes non visitées.
Un snapshot est généré à chaque extension du circuit. Complexité : $O(E)$.

% ─────────────────────────────────────────────────────────────────────────
\section*{Workflow interface $\rightarrow$ algorithmes}
% ─────────────────────────────────────────────────────────────────────────
 

Le workflow de données suit un flux unidirectionnel :

\begin{enumerate}
    \item L'utilisateur crée ou modifie le graphe via l'éditeur visuel ou la matrice.
    \item Chaque action est dispatchée au reducer qui met à jour l'état \texttt{GraphUI}.
    \item L'utilisateur sélectionne un algorithme dans le menu (ou saisit une requête en canvas).
    \item Le module \texttt{algorithmCinema} reçoit l'objet \texttt{GraphUI} et produit une séquence
          de snapshots (états intermédiaires).
    \item Le module \textit{Algorithm Cinema} rejoue ces snapshots dans le canvas SVG,
          permettant une visualisation pas à pas avec contrôles play/pause/step/rewind/scrub.
\end{enumerate}

\bigskip
 
 
% ─────────────────────────────────────────────────────────────────────────
\section{Tests et validation}
% ─────────────────────────────────────────────────────────────────────────
 
\subsection{Tests nominaux — Labibi Idris et Drira Yassine}
 
Les algorithmes ont été validés sur des structures calculables à la main. Le tableau
suivant résume les cas de test nominaux utilisés, avec des exemples concrets de graphes.

\renewcommand{\arraystretch}{1.3}
\begin{longtable}{>{\bfseries}p{3.2cm} p{4.8cm} p{5.2cm}}
\rowcolor{headerblue}
\color{white}\textbf{Algorithme} & \color{white}\textbf{Graphe de test} & \color{white}\textbf{Résultat attendu et vérifié} \\
\hline
\endfirsthead
\rowcolor{headerblue}
\color{white}\textbf{Algorithme} & \color{white}\textbf{Graphe de test} & \color{white}\textbf{Résultat attendu et vérifié} \\
\hline
\endhead

\rowcolor{teal100}
Dijkstra &
Chemin linéaire $P_5$ : $1\xrightarrow{2}2\xrightarrow{3}3\xrightarrow{1}4\xrightarrow{4}5$, source = 1 &
Distances : $d(1,2)=2$, $d(1,3)=5$, $d(1,4)=6$, $d(1,5)=10$. \\

Dijkstra &
$K_4$ avec poids : arêtes $(1,2)=1$, $(1,3)=4$, $(1,4)=3$, $(2,3)=2$, $(2,4)=5$, $(3,4)=1)$ ; source = 1 &
Distances : $d(2)=1$, $d(3)=3$, $d(4)=3$. Chemin optimal $1\to3$ passe par $1\to2\to3$. \\

\rowcolor{teal100}
Bellman-Ford &
Graphe orienté à 4 sommets : $(1,2,-1)$, $(1,3,4)$, $(2,3,3)$, $(2,4,2)$, $(3,4,5)$, source = 1 &
Distances : $d(2)=-1$, $d(3)=2$, $d(4)=1$ (chemin $1\to2\to4$). \\

Prim / Kruskal &
Graphe non orienté $K_4$ pondéré : arêtes $(1,2)=4$, $(1,3)=2$, $(2,3)=1$, $(2,4)=5$, $(3,4)=8)$ &
ACM = $\{(2,3), (1,3), (2,4)\}$, poids total = 8. \\

\rowcolor{teal100}
Kruskal &
Cycle $C_5$ : 5 sommets en anneau, poids $1,2,3,4,5$ &
ACM : 4 arêtes acceptées (poids les plus faibles), l'arête de poids 5 rejetée (cycle). \\

Edmonds-Karp &
Réseau de flot à 4 sommets : source $s=1$, puits $t=4$, capacités $(1,2)=3$, $(1,3)=2$, $(2,4)=2$, $(3,4)=3$, $(2,3)=1)$ &
Flot maximum = 4 (deux chemins augmentants : $s\to1\to4$ et $s\to2\to4$). \\

\rowcolor{teal100}
Composantes connexes &
Graphe à 6 sommets, deux composantes : $\{1,2,3\}$ en triangle et $\{4,5,6\}$ en triangle &
Deux composantes identifiées, colorées distinctement. \\

CFC (Kosaraju / Tarjan) &
Graphe orienté : $1\to2\to3\to1$ (cycle) et $3\to4$ (pont vers un sommet isolé) &
Deux CFC : $\{1,2,3\}$ et $\{4\}$. \\

\rowcolor{teal100}
Greedy coloring &
Cycle $C_4$ (graphe biparti) : $1-2-3-4-1$ &
2 couleurs suffisent : $\{1,3\}$ couleur A, $\{2,4\}$ couleur B. \\

Greedy coloring &
$K_4$ complet (4 sommets tous reliés) &
4 couleurs nécessaires, une par sommet. \\

\rowcolor{teal100}
Hierholzer &
Graphe non orienté : cycle $C_4$ (tous degrés = 2) &
Circuit eulérien trouvé : $1\to2\to3\to4\to1$. \\

Hierholzer &
Graphe orienté à 3 sommets : $1\to2\to3\to1$ ($\deg^+=\deg^-=1$ pour chacun) &
Circuit eulérien orienté : $1\to2\to3\to1$. \\

\end{longtable}

Des graphes plus grands générés aléatoirement ($n = 20$--$50$ sommets) ont également
été utilisés pour évaluer les performances et la robustesse de la visualisation.
 
\subsection{Tests des cas limites — Fathallah Youssef et Rayen Mestiri}
 
\renewcommand{\arraystretch}{1.3}
\begin{longtable}{>{\bfseries}p{4cm} p{9.5cm}}
\rowcolor{headerblue}
\color{white}\textbf{Cas limite} & \color{white}\textbf{Comportement attendu et vérifié} \\
\hline
\endfirsthead
\rowcolor{headerblue}
\color{white}\textbf{Cas limite} & \color{white}\textbf{Comportement attendu et vérifié} \\
\hline
\endhead
 
\rowcolor{teal100}
Graphe vide ($|V|=0$) &
Tous les algorithmes retournent un résultat vide ou affichent un message explicatif sans lever d'exception. \\
 
Graphe à 1 sommet &
Dijkstra et Bellman-Ford retournent une distance nulle pour la source ; Prim et Kruskal retournent un arbre vide ; Hierholzer confirme l'absence de circuit. \\
 
\rowcolor{teal100}
Graphe non connexe &
Dijkstra et Prim assignent $+\infty$ aux sommets inaccessibles et le signalent dans le panneau de résultats. Les composantes connexes identifient correctement chaque sous-graphe isolé. \\
 
Poids négatifs (Dijkstra) &
Dijkstra refuse l'exécution et affiche un message redirigeant vers Bellman-Ford. Bellman-Ford traite le graphe correctement, y compris avec plusieurs arêtes de poids négatif. \\
 
\rowcolor{teal100}
Cycle négatif &
Ex. : graphe orienté $1\xrightarrow{1}2\xrightarrow{-3}3\xrightarrow{1}1$ (cycle de poids $-1$). Bellman-Ford détecte le cycle lors de la $|V|$-ème itération et affiche une alerte avec l'arête incriminée. \\
 
Graphe non eulérien (non orienté) &
Ex. : graphe en étoile $K_{1,3}$ (le centre a degré 3). Hierholzer liste les sommets de degré impair et stoppe l'exécution avec un message explicatif. \\
 
\rowcolor{teal100}
Graphe non eulérien (orienté) &
Ex. : graphe $1\to2$, $1\to3$, $2\to3$ ($\deg^+(1)=2 \neq \deg^-(1)=0$). Les sommets déséquilibrés sont signalés. \\
 
Boucles et multi-arêtes &
Les structures Union-Find et les listes d'adjacence gèrent correctement les arêtes multiples et les boucles sans corruption de l'état du graphe. \\
 
\end{longtable}
 
% ─────────────────────────────────────────────────────────────────────────
\section{Lancement de la plateforme}
% ─────────────────────────────────────────────────────────────────────────
 
\begin{verbatim}
git clone https://github.com/tasnime-bbker/GraphLab
cd GraphLab
npm install
npm run dev        # http://localhost:5173  (Vite dev server)
npm run build      # build de production (tsc + vite build)
\end{verbatim}
 
\clearpage


\clearpage


% ─── CONCLUSION ──────────────────────────────────────────────────────────────
\chapter*{Conclusion générale}
\addcontentsline{toc}{chapter}{Conclusion générale}

\section*{Bilan}

Le projet Graph Lab a abouti à la livraison d'une plateforme fonctionnelle et interactive
permettant de visualiser et comprendre les principaux algorithmes de graphes. L'ensemble
des fonctionnalités planifiées a été implémenté, testé et intégré dans un produit cohérent.
La collaboration entre les onze membres de l'équipe, organisée autour d'un dépôt GitHub
commun et de daily standups réguliers, a permis de maintenir un rythme de développement
soutenu tout en limitant les conflits d'intégration. La transversalité des équipes front
et test a notamment renforcé la qualité du produit final.

\section*{Difficultés rencontrées}

\begin{itemize}
    \item La synchronisation entre l'éditeur visuel SVG et la saisie matricielle a nécessité
          une conception rigoureuse du reducer centralisé pour éviter les incohérences d'état.
    \item La génération des snapshots de visualisation a dû être coordonnée entre les développeurs
          algorithmes et l'équipe front, créant une dépendance inter-modules à gérer prudemment.
    \item La détection correcte des cycles négatifs dans Bellman-Ford a requis une attention
          particulière à la condition de terminaison et à l'identification précise de l'arête coupable.
\end{itemize}

\section*{Améliorations futures}

\begin{itemize}
    \item Mode collaboratif en temps réel (plusieurs utilisateurs sur le même graphe).
    \item Ajout d'un mode quiz pédagogique permettant à l'utilisateur de prédire la prochaine étape.
\end{itemize}

\end{document}